\documentclass[11pt]{article}
\usepackage[margin=1in]{geometry}
\usepackage{amsmath,amssymb,amsfonts,mathtools,bm}
\usepackage[T1]{fontenc}
\usepackage[utf8]{inputenc}
\title{Inline and Display Mathematics}
\author{Source-backed grouped formula sample}
\date{}
\begin{document}
\maketitle
\section{Expressions}
\paragraph{Expression 1.} The following expression is taken from a source-backed formula corpus.

\begin{align*}\mathcal{I}_{\mathrm{D} }(\delta) = \sup_{\gamma \in \Sigma_{\mathrm{D} }(\delta)} \int_{\mathcal{Y}_1 \times \mathcal{Y}_2} f(y_1, y_2) \, d \gamma(y_1, y_2) \mathcal{I}(\delta) = \sup_{\gamma \in \Sigma(\delta)} \int_{\mathcal{S}} f(y_1, y_2) \, d \gamma(y_1, y_2, x).\end{align*}

\paragraph{Expression 2.} The following expression is taken from a source-backed formula corpus.

\begin{align*}\mathcal{I}(\delta_1, 0) = \widehat{\mathcal{I}} (\delta_1, \delta_2) = \inf _{\lambda_1 \in \mathbb{R}_{+}}\left[\langle\lambda, \delta\rangle +\sup _{\varpi \in \Pi\left(\mu_{13}, \mu_{23}\right)} \int_{\mathcal{V}} f_{\lambda}(s_1, s_2) \, d \varpi(s_1,s_2) \right] ,\end{align*}

\paragraph{Expression 3.} The following expression is taken from a source-backed formula corpus.

\begin{align*}y^\delta \tilde{\Gamma}^\alpha {}_{\beta \delta }\tilde{B}^\beta|_{1+3}=y^\delta dx^\sigma [\Gamma ^\alpha {}_{\beta \delta }B^\beta{}_{,\sigma }+\Gamma ^\alpha {}_{\beta \delta }\Gamma ^\beta {}_{\sigma \tau}B^\tau +\Gamma ^\alpha {}_{\beta \delta ,\sigma }B^\beta ]. \end{align*}

\paragraph{Expression 4.} The following expression is taken from a source-backed formula corpus.

\begin{align*}\theta _{\bot \parallel }=\frac{1}{y^{2}}\left( \partial _{\tau }\vec{x}\, \partial _{\sigma }\vec{x}+\partial _{\tau }y\, \partial _{\sigma }y\right) =\frac{1}{a^{2}\tau }\left[ \vec{f}{\vec{c}\, }'+aa'\right] +\frac{1}{a^{2}}\left( \vec{g}{\vec{c}\, }'\right) +\ldots\end{align*}

\paragraph{Expression 5.} The following expression is taken from a source-backed formula corpus.

\begin{align*}\left. \frac{d}{d\tau}\right|_{\tau=0} x^{\alpha}(g_{A}(\tau)g)=\left.\frac{\partial x^{\alpha}(g_{A}(\tau)g)}{\partial x^{M}(g_{A}(\tau))}\right|_{g_{A}=e}\cdot\left.\frac{d x^{M}(g_{A}(\tau))}{d \tau}\right|_{\tau=0}=R^{\alpha}_{M}(g)\delta^{M}_{A}=R^{\alpha}_{A}(g).\end{align*}
\end{document}
